\originalchapter{辐射耦合与跃迁强度}
\originalsection{最小耦合与电偶极表象}
电荷$q$的最小耦合为$[\mathbf p-q\mathbf A]^2/(2m)+q\phi$. Coulomb规范$\nabla\cdot\mathbf A=0$下, 交互作用包括$-q\mathbf p\cdot\mathbf A/m$与$q^2A^2/(2m)$. 在原子大小远小于光波长的长波近似中, 经相应规范变换可用$-\mathbf d\cdot\mathbf E$组织同一物理过程. 不能在一个表象中截掉必要的$A^2$项, 再与另一表象的完整结果比较.

对本征态, $[H_0,\mathbf r]=-\ii\hbar\mathbf p/m$, 所以$\bra b\mathbf p\ket a=\ii m\omega_{ba}\bra b\mathbf r\ket a$. 完整中间态求和与此关系保证长度、速度表象一致. 这在Rayleigh与Thomson散射中尤其重要: 低频响应主要由虚激发的极化描述, 高频响应趋近自由电子, 两个表象的方便程度不同.

\originalsection{选择定则与振子强度}
电偶极是奇宇称的一级球张量. 按角动量耦合, $\Delta m=q=0,\pm1$, $|j_b-j_a|\le1\le j_b+j_a$, 因而$0\leftrightarrow0$禁戒; 轨道部分$\Delta l=\pm1$, 忽略自旋轨道混合时$\Delta s=0$. 写成Wigner--Eckart形式,
\[\bra{j_bm_b}d_q\ket{j_am_a}=(-1)^{j_b-m_b}\begin{pmatrix}j_b&1&j_a\\-m_b&q&m_a\end{pmatrix}\bra{j_b}\!\!\left|\mathbf d\right|\!\!\ket{j_a}.\]
同一对角动量多重态的相对线强由$3j$系数平方决定, 径向结构归入约化矩阵元. 这是计算偏振分支比而不用重复径向积分的理由.

对一个Cartesian方向定义$f_{ba}=2m\omega_{ba}|\bra b x\ket a|^2/\hbar$. 由$\bra a[x,[H_0,x]]\ket a=\hbar^2/m$并插入完备态, 左边等于$2\sum_b(E_b-E_a)|x_{ba}|^2$, 因而$\sum_b f_{ba}=1$. 求和必须包括连续谱; 激发态向下跃迁的$f$为负. 多电子总和为电子数, 不能对每一条共振都设$f=1$.

\originalsection{Einstein系数与真空自发辐射}
每单位体积、每单位角频率的三维电磁模密度为$\omega^2/(\pi^2c^3)$, 已计两种偏振. 热平均每模占据$\bar n=(\e^{\hbar\omega/k_BT}-1)^{-1}$, 因而$u_\omega=\hbar\omega^3/[\pi^2c^3(\e^{\hbar\omega/k_BT}-1)]$. 两级统计权重$g_g,g_e$时, 详细平衡$N_gB_{ge}u=N_e(A_{eg}+B_{eg}u)$和Boltzmann比给$g_gB_{ge}=g_eB_{eg}$,$A_{eg}/B_{eg}=\hbar\omega^3/(\pi^2c^3)$.

量子场的产生算符有$\bra{n+1}a^\dagger\ket n=\sqrt{n+1}$, 所以发射速率含$n+1$, 吸收含$n$. 其中$1$是自发辐射, 不是假定真空有经典随机光子的结果. 将每模电场$\sqrt{\hbar\omega/(2\epsilon_0V)}$代入黄金律, 对传播方向与偏振积分, 用$\int\dd\Omega[|\mathbf d|^2-|\mathbf d\cdot\hat k|^2]=8\pi|\mathbf d|^2/3$, 得
\begin{equation}\Gamma_{eg}=\frac{\omega_{eg}^3|\mathbf d_{eg}|^2}{3\pi\epsilon_0\hbar c^3}.\label{eq:spont}\end{equation}
这要求偶极近似、连续宽带模、弱耦合和近似平坦的局域态密度. 高品质腔中单模的可逆交换不能直接视为指数衰减.
